\section{Large Scale Structure Formation}

Now that we have understood how modern cosmology came to be, we can move on to understanding the formation of structures. Observations from galaxy surveys such as DESI show that 
although on large scales the Universe looks homogenous, the smaller scales indicate the existence of an underlying non-homogenous structure. This can be seen through the formation of galaxy clusters 
and the cosmic web. Understanding how these structures are formed gives us a lot of information regarding the strength of gravity at large scales as well as the amount of matter in the Universe.
Since the full derivation of each equation in structure formation is too long to write for this thesis, we recommend reading \cite{Knobel_2013} and \cite{Dodelson_2003} for a more detailed derivation\\


\subsection{Newtonian approach to structure formation}

In this introductory section, we will be starting off by describing structure formation by going only to linear scales. This will be the model on which we build up more complex 
structure formation model. \\

\noindent \textbf{\underline{Theory of structure formation}} \\

We start by assuming a Universe is filled with an inhomogeneous, dissipationless, ideal fluid with matter density $\rho(\mathbf{r},t)$, velocity field $\vec{v}(\mathbf{r}, t)$, 
pressure $P(\mathbf{r},t)$ and entropy per unit mass $S(\mathbf{r}, t)$ with $\mathbf{r}$ being the fluid's coordinates. We obtain the following hydrodynamical equations:
\begin{equation}
    \begin{cases}
        \frac{\partial \rho}{\partial t} + \vec{\nabla}. (\rho \vec{v}) = 0 & \mathrm{Continuity\; equation\;} \\
        \frac{\partial \vec{v}}{\partial t} + (\vec{v} . \vec{\nabla}) \vec{v} + \frac{\vec{\nabla} P }{\rho} + \vec{\nabla}\Phi = 0 & \mathrm{Euler\; equation\;} \\
        \nabla^2 \Phi = 4 \pi G \rho & \mathrm{Poisson\; equation\;} \\
        \frac{\partial S}{\partial t } + (\vec{v} . \vec{\nabla}) S = 0& \mathrm{Entropy\; conservation\;}
    \end{cases}
    \label{eqn:hydro_eqn}
\end{equation}
with $\Phi$ the gravitational potential. Since matter is approximated as a single-stream fluid, 
neighboring particles follow nearly parallel trajectories without "shell crossing" (collision of streams). This description breaks down once nonlinear collapse occurs, 
when multiple streams overlap and complex dynamics emerge. A fully accurate treatment of structure formation would instead 
require solving the general relativistic Boltzmann equation including all cosmic components and their interactions, although this is in itself a difficult problem due to the nonlinear 
coupling between the fluid equations. To overcome that issue, we assume a FLRW Universe and perturb the fluid around the Hubble flow, meaning we 
separate each variable as a homogenous part and an inhomogeneous perturbation:
\begin{align*} 
    \rho(\mathbf{r},t) &=  \bar{\rho}(t) + \delta \rho(\mathbf{r},t) \\ 
    \mathbf{v}(\mathbf{r},t) &=  \bar{\mathbf{v}}(\mathbf{r},t) + \delta \mathbf{v}(\mathbf{r},t) \\
    P(\mathbf{r},t) &=  \bar{P}(t) + \delta P(\mathbf{r},t) \\
    \Phi(\mathbf{r},t) &=  \bar{\Phi}(\mathbf{r},t) + \delta \Phi(\mathbf{r},t) \\
    S(\mathbf{r},t) &=  \bar{S}(t) + \delta S(\mathbf{r},t) \\
\end{align*}
\textcolor{red}{COMMENT : detail why not relativistic treatment and small perturbations}

We can then re-introduce our new variables in Equations \ref{eqn:hydro_eqn} although to slightly simplify the equations we will be working in comoving coordinates $\mathbf{x}$ such that 
\begin{align*}
    \mathbf{r} = a \mathbf{x} \\
    \nabla_{\mathbf{r}} = \frac{1}{a} \nabla_{\mathbf{x}} \\
    \frac{\partial}{\partial t}\bigg\rvert_{\mathbf{r}} = \frac{\partial}{\partial t}\bigg\rvert_{\mathbf{x}} - \frac{1}{a}\mathbf{v} . \nabla_{\mathbf{x}} \\
\end{align*}
We will also be working in Fourrier space where we get for any perturbation $\delta q(\mathbf{x}, t)$:
\begin{equation}
    \delta q(t, \mathbf{k}) = \int \delta q (t, \mathbf{x}) e^{-i \mathbf{k}\mathbf{x}} \mathrm{d}^3x
\end{equation}

\begin{equation}
    \delta q(t, \mathbf{x}) = \frac{1}{(2\pi)^3} \int \delta q (t, \mathbf{k}) e^{-i \mathbf{k}\mathbf{x}} \mathrm{d}^3k
\end{equation}
With $\mathbf{k}$ are the comoving Fourrier modes. We get the following hydrodynamical equations at first order:
\begin{equation}
    \boxed{
    \begin{cases}
        \delta \dot{\rho} + 3H\delta \rho + \frac{i\bar{\rho}\mathbf{k}}{a} . \delta\mathbf{v} = 0& \mathrm{Continuity\; equation\;} \\
        \delta \dot{\mathbf{v}} + H \delta \mathbf{v} + \frac{i\mathbf{k}}{a\bar{\rho}}(c_s^2 \delta \rho + \sigma \delta S) + \frac{i \mathbf{k}}{a} \delta \Phi = 0 & \mathrm{Euler\; equation\;} \\
        k^2 \delta \Phi + 4 \pi G a^2 \delta \rho = 0 & \mathrm{Poisson\; equation\;} \\
        \delta \dot{S} = 0& \mathrm{Entropy\; conservation\;}
    \end{cases}
    }
    \label{eqn:newton_hydro_eqn}
\end{equation}
with $c_s^2 = (\partial \bar{P}/ \partial \bar{\rho})_{\bar{S}}$ the speed of sound in the fluid, $\sigma = (\partial \bar{P}/ \partial \bar{S})_{\bar{\rho}}$ and $k=|\mathbf{k}|$. \\

\textcolor{red}{COMMENT : how do we get the 3h delta rho?} \\


What we are interested in is essentially how do matter perturbation grows as they will give birth to structure by gravitational attraction. The first thing to note is 
that the entropy fluctuations must exist at early times in order for them to keep existing otherwise they would simply dissipate. Moreover, results from the simplest 
models of inflation do not support the fluctuations of entropy. So from now on we well set $\delta S = 0$. We can also assume that the velocity perturbations would follow the overdense regions giving us 
$\mathbf{k} || \delta \mathbf{v}$. We call the modes that result from these conditions the adiabatic modes. \\
\textcolor{red}{COMMENT : say no vorticity to justify parallel and separate from adiabatic modes}.

We will be denoting the overdensity $\delta(\mathbf{k},t)$ as the matter density contrast which is written as:
\begin{equation}
    \delta(\mathbf{k}, t) = \frac{\rho(\mathbf{k},t)}{\bar{\rho}(t)} - 1 = \frac{\delta \rho (\mathbf{k}, t)}{\bar{\rho}(t)}
\end{equation}

In the adiabatic modes, the continuity equation becomes:
\begin{equation}
    \dot{\delta} + \frac{ik}{a}\delta v = 0
\end{equation}
If we then differentiate this equation and combining its results with Euler and Poisson equations we obtain:
\begin{equation}
    \ddot{\delta} + 2 H \dot{\delta} + \left( \frac{c_s^2 k^2}{a^2} - 4\pi G \bar{\rho} \right)\delta = 0
    \label{eqn:linear_dens}
\end{equation}
Here we obtain a linear second order differential equation with two solutions to it. We denote $\lambda = \frac{2\pi a}{k}$ the physical wavelength of the mode, 
$\lambda_J = c_s \sqrt{\frac{\pi}{G\rho}} $ the Jeans length and $M_J = \frac{\pi}{6} \lambda_J^3 \bar{\rho} $ the Jeans mass. Equation \ref{eqn:linear_dens} tells us that 
if $\lambda < \lambda_J$, the solution to the equation is an oscillating wave leaving no room for structures to grow. So a condition we must have is $\lambda > \lambda_J$. We 
place ourselves in the condition where the wavelength is much larger than the Jeans length $\lambda \gg \lambda_J$, Equation \ref{eqn:linear_dens} is now reduced to:
\begin{equation}
    \ddot{\delta} + 2 H \dot{\delta} - 4\pi G \bar{\rho} \delta = 0
\end{equation}

\textcolor{red}{COMMENT : important to specify this reduction for collisionless matter so DM but not baryonic} \\


which has the following general solution:
\begin{equation}
    \delta(\mathbf{k},t) = \delta_+(\mathbf{k})D_+(t) + \delta_-(\mathbf{k})D_-(t)
\end{equation}
with the + and - denoting growing and decaying modes. We call the $D_+(t)$ term the \textbf{linear growth function}. In the radiation dominated era, the Mészáros 
effect suppresses the growth of sub-Hubble matter perturbations, which grow only logarithmically rather than as a power law. In the matter dominated era 
the expansion rate grows as $H(t) = \frac{2}{3t}$ which gives us:
\begin{equation}
    \boxed{D_+ \propto t^{2/3} \propto a(t)}
\end{equation}
\begin{equation}
    D_- \propto t^{-1}
\end{equation}
This means that we can get grown modes for matter perturbations in the matter dominated era with $D_+$ while $D_-$ goes to 0. A general formula for $D_+$ can be obtained 
(\cite{Carroll_1992}) and it gives us:
\begin{equation}
    D_+(z) = \frac{1}{1+z} \frac{g(z)}{g(0)}
\end{equation}
\begin{equation}
    g(z) = \frac{\Omega_m(z)}{\Omega^{4/7}_m(z)-\Omega_{\Lambda}(z) + (1+\Omega_m(z)/2)(1+\Omega_{\Lambda}(z)/70)}
\end{equation}

The growth of perturbations described by $D_+$ can be characterised by the \textit{growth rate} $f$, defined as the logarithmic derivative of the growth factor 
with respect to the scale factor:
\begin{equation}
    f = \frac{d \ln D_+}{d \ln a}
\end{equation}
In the matter dominated era, since $D_+ \propto a$, we get $f = 1$. A useful approximation valid across different epochs and cosmologies is given by 
(\cite{Peebles_1980}):
\begin{equation}
    \boxed{f \approx \Omega_m(z)^{\gamma}}
\end{equation}
where $\gamma$ is obtained from the underlying gravitational theory. For a flat $\Lambda$CDM cosmology and GR it is $\gamma \simeq 0.55$. 
Deviations from this value would indicate modifications to GR on cosmological scales. 
In practice, $f$ is not directly observable, and is instead measured in combination with $\sigma_8$, the root mean square of matter fluctuations on scales of 
$8 \, h^{-1}$Mpc. The combination $f\sigma_8$ is therefore the key observable quantity that encodes 
the rate of growth of large-scale structure and provides a direct test of gravity on cosmological scales. \\

\textcolor{red}{COMMENT : not just Peebles more ref and there is an expression for gamma wrt w} \\
\textcolor{red}{COMMENT : can relate here or later As conneciton} \\


\noindent \textbf{\underline{Statistical analysis of overdensities}} \\

Practically, we perform a statistical analysis when it comes to the study of the structures of the Universe. $\delta$ is looked at as a realization of a stochastic process 
and by analyzing that we can compare our observations to theories. \\

\textcolor{red}{COMMENT : primordial fluctuations expected from inflation described in later section implies stochastic nature etc... can relate to CMB observations 
today (just to make the link)} \\

We introduce $\langle...\rangle$ as the ensemble average of the stochastic process underlying a random field. The large-scale structure we observe is only one realization 
of an underlying random cosmic process, so estimating its true statistical properties requires additional assumptions. A common assumption is ergodicity, meaning that 
averages measured over a sufficiently large volume of the Universe approximate the averages over many hypothetical realizations. In practice, this is related to the fair 
sample hypothesis, which assumes that widely separated regions of the Universe can be treated as nearly independent samples of the same process. If the survey volume is 
large enough and correlations become weak on large scales, these approximations are reasonable. However, real galaxy surveys cover finite volumes and therefore their measured 
averages are affected by statistical fluctuations known as \textbf{sample variance}, or \textbf{cosmic variance} when limited by the size of the observable Universe.\\

We denote the 2-point correlation function as:
\begin{equation}
    \xi(\mathbf{x}, \mathbf{x}') = \left\langle \delta(\mathbf{x}) \delta(\mathbf{x}') \right\rangle
\end{equation}
Our field is homogenous and isotropic therefore the only 
quantity that affects the correlation function is $r=|\mathbf{x} - \mathbf{x}'|$, $\xi(\mathbf{x}, \mathbf{x}')=\xi(r)$. If at $r=+\inf$, the correlations go to 0, 
then we can decompose it in Fourrier modes:
$\delta(\mathbf{k})$:
\begin{equation}
    \left\langle \delta(\mathbf{k}) \delta^*(\mathbf{k}') \right\rangle = (2 \pi)^3 \delta_D(\mathbf{k} - \mathbf{k}')\mathcal{P}(\mathbf{k})
\end{equation}
\begin{equation}
    \mathcal{P}(\mathbf{k}) = \int \xi(\mathbf{x})e^{-i \mathbf{k} \mathbf{x}} d^3x
\end{equation}
where the $\mathcal{P}(k)$ is the \textbf{power spectrum}, which is a real function, and $\delta_D$ is the Dirac delta function. \\
The works from \cite{Harrison_1970}, \cite{Zeldovich_1970_pk} and \cite{Peebles_1970} show that there is a preferrable form for the initial primordial power spectrum that takes the form:
\begin{equation}
    \mathcal{P}(k) \propto k^{n_s}
\end{equation}
with $n_s$ the \textbf{spectral index}. If $n_s=1$ the power spectrum becomes scale invariant. It is possible to predict the shape of the power spectrum of the initial conditions by studying an inflationary model. Most common models 
and results from \cite{Planck_2018} indicate that the initial perturbations are Gaussian with $n_s \approx 1 $. \\
The power spectrum for the linear overdensity field is:
\begin{equation}
    \mathcal{P}_{\mathrm{lin}}(t, k) = A_s k^{n_s}T^2(k)D^2_+(t)
\end{equation}
with $A_s$ the normalization parameter and $T(k)$ is the transfer function which characterizes how each Fourrier mode of the density field evolved from the initial conditions to 
a later time. \\
\textcolor{red}{COMMENT : can link sigma8 and As here} \\


\subsection{Relativistic structure formation}

Similarly to the previous section, we will be introducing structure formation from a general relativistic framework with a focus on linear perturbation. We will not be deriving 
every single quantity, however explaination of how we get there and how it relates to the Newtonian approximation will be present. \\

\noindent \textbf{\underline{General description of perturbations}} \\

Firstly, we will be working in the FLRW metric for $k=0$. We will also work in conformal time $\tau$ which simplifies this problem. 
It is defined as $a \mathrm{d}\tau = \mathrm{d}t$. We will also work in natural units setting $c=\hbar =1$. We will denote the FLRW metric as $\bar{g}_{\mu \nu}$ so that 
for a flat spacetime we have:
\begin{equation}
    ds^2 = \bar{g}_{\mu \nu} dx^{\mu} dx^{\nu} = a^2(\tau)\left( - d\tau^2 + \bar{\gamma}_{ij} dx^{i} dx^{j} \right)
\end{equation}
where $\bar{\gamma}_{ij}$ is the metric for the spatial components and $(i,j)$ can take the values $1,2,3$ (in general we will put a bar on top of the parameter related to a flat FLRW metric). \\

In a relativistic framework, the perturbations are done on the metric in itself such that when we choose a metric $g_{\mu \nu}$ the deviations from the homogenous isotropic 
background are close enough to $\bar{g}_{\mu \nu}$\footnote{Mathematically, this change in metric means we alsothe mainifolds being worked on so simple mathematical operations
no longer apply. One way to overcome that issue is to define $\bar{g}_{\mu \nu}$'s functional form as invariant under coordinate transformations.}. We then obtain:
\begin{equation}
    \delta g_{\mu \nu} = g_{\mu \nu}(x) - \bar{g}_{\mu \nu}(x)
\end{equation} 
Physically, we are splitting $g_{\mu \nu}$ into a background $\bar{g}_{\mu \nu}$ and perturbation $ \delta g_{\mu \nu}$. However this perturbation is not a 4-tensor which will 
require us to set a gauge transformation. Similarly, the perturbation for a general scalar-like object and a vecotr-like objects are defined as:
\begin{equation}
    \delta q(x) = q(x) - \bar{q}(x)
\end{equation}
\begin{equation}
    \delta q^{\mu} = q^{\mu} - \bar{q}^{\mu} = \delta q_{\nu} \bar{g}^{\mu \nu} + \bar{q}_{\nu} \delta g^{\mu \nu}
\end{equation}
We can also perform the same for the energy-momentum tensor.
\begin{equation}
    \delta T_{\mu \nu} = T_{\mu \nu} - \bar{T}_{\mu \nu}
\end{equation}
with 
\begin{equation}
    T_{\mu \nu} = (\rho + P)u_{\mu}u_{\nu} + P g_{\mu \nu} + \Pi_{\mu \nu}
\end{equation}
and $\Pi_{\mu \nu}$ is the anisotropic stress which accounts for deviations from an ideal fluid. \\

The general form for the metric perturbation is written as:
\begin{equation}
    \delta g_{\mu \nu} dx^{\mu} dx^{\nu} = a^2(\tau)\left[ -2A d\tau^2 + B_i d\tau dx^i + 2(H_L \delta_{ij} + H_{ij})dx^i dx^j \right]
\end{equation}
where $A(\tau, \mathbf{x})$, $B_i(\tau, \mathbf{x})$, $H_L(\tau, \mathbf{x})$ and $H_{ij}(\tau, \mathbf{x})$ are perturbation quantities and $\delta_{ij}$ a delta function.\\

\noindent \textbf{\underline{Newtonian gauge}} \\

These perturbation quantities intervene will then intervene in the field equations. Determining them necessitates a gauge transformation. We will be showing the results 
of the Newtonian gauge as it is useful for analyzing the density fluctuations inside of the Hubble horizon ($D_H = 1/H$) as well as it is simple to interpret its results in terms of 
Newtonian physics. The gauge is defined by:
\begin{equation}
    B=H=0
\end{equation}
and we will be renaiming the other perturbations as:
\begin{equation}
    \Psi = A
\end{equation}
\begin{equation}
    \Phi = -H_L
\end{equation}

The field equations for the Newtonian gauge are:
\begin{equation}
    \boxed{
    \begin{cases}
        -k^2 \Phi - 3H(\dot{\Phi} + H\Psi) = 4 \pi G a^2 \delta \rho \\
        \dot{\Phi} + H \Psi = 4 \pi G a^2 \frac{v}{k} (\bar{\rho} + \bar{P}) \\
        \ddot{\Phi} + H(\dot{\Psi} + 2\dot{\Phi}) + (2 \dot{H} + H^2)\Psi + \frac{k^2}{3}(\Phi - \Psi) = 4 \pi G a^2 \delta P \\
        k^2 (\Phi - \Psi) = 8 \pi G a^2 \Pi
    \end{cases}
    }
\end{equation}
where the dots are derivatives with respect to the conformal time. 
If we consider scale much smaller than the Hubble horizon $H/k \ll 1$, the first field equation then becomes:
\begin{equation}
    -k^2 \Phi = 4 \pi G a^2 \delta \rho
\end{equation}
In this limit we recover the Poisson equation in the Newtonian formalism. Therefore, we can say that well within the horizon the Newtonian gauge reduces to Newtonian mechanics. 
This also means we can have a physical interpretation for $\Phi$ it acting as the Newtonian gravitational potential (called generalized Newtonian potential). The last equation 
also informs us that in the matter dominated epoch we have $\Phi \simeq \Psi$ reducing the metric perturbations to only one parameter. \\

One can only also derive the equations of motion in the Newtonian gauge, however for this thesis we will only be showing the final equations of motions obtained for matter 
fluctuations in the matter dominated era. We note the matter overdensity $\delta_m = \frac{\delta \rho_m}{\bar{\rho}}$, we obtain:
\begin{equation}
    \ddot{\delta}_m + H \dot{\delta}_m - 4\pi G a^2 \bar{\rho}_m\delta_m = 0
\end{equation}

\subsection{Inflation}

Inflation is a period of time at the very early stages of the Universe 
where an accelerated expansion of the Universe occurred (\cite{Guth_1981}). It was theorized to tackle known issues regarding the fine tuning of the flatness of the Universe and the horizon problem.
What we see as well is that through inflation we can have a predict for the initial density fluctuations. 
This following condition is imposed on inflation to obtain an accelerated early Universe:
\begin{equation}
    \ddot{a} > 0
\end{equation}

During inflation, we assume that the Universe is dominated by a quantum scalar field $\phi$ called the inflaton, experiencing a "slow roll" down its potential $V(\phi)$ (\cite{Mukhanov_1992}, \cite{Albrecht_1982}). 
The action for a scalar field minimally coupled to gravity is given by:
\begin{equation}
    S = \int d^4x \sqrt{-g} \left[ \frac{M^2_{\mathrm{pl}}}{2} R - \frac{1}{2} g^{\mu \nu} \partial_{\mu}\phi \partial_{\nu}\phi - V(\phi) \right]
\end{equation}
where $M_{\mathrm{pl}} $ is the Planck mass defined as $M^2_{\mathrm{pl}} = \left(\frac{1}{8\pi G}\right)^2$. Varying the action with the respect to the metric allows to obtain 
the energy momentum tensor $T_{\mu \nu}$ and its components $P_{\phi}$ and $\rho_{\phi}$:
\begin{equation}
    P_{\phi} = \frac{1}{2} \dot{\phi}^2 - V(\phi)
\end{equation}
\begin{equation}
    \rho_{\phi} = \frac{1}{2} \dot{\phi}^2 + V(\phi)
\end{equation}
where the kinetic energy of the field is given by $\frac{1}{2} \dot{\phi}^2$.\\
The dynamics can be recovered with the Klein Gordon equation giving us:
\begin{equation}
    \ddot{\phi} + 3H \dot{\phi} + \frac{d V(\phi)}{d \phi} = 0
\end{equation}

We can also apply the Friedmann equations to the inflaton field where we get:
\begin{equation}
    H^2 = \frac{1}{M^2_{\mathrm{pl}}}\left( \frac{1}{2} \dot{\phi}^2 + V(\phi) \right)
\end{equation}
\begin{equation}
    \dot{H} = -\frac{1}{2} \frac{\dot{\phi}^2}{3M^2_{\mathrm{pl}}}
\end{equation}

Assuming a slow rolling field entails that kinetic term is small compared to the potential term $V$. We can show that given this assumption $a \sim e^{Ht}$ which is a necessary
condition to inflation. Additionally we assume coherent oscillations around a minimum that are damped by a $H\dot{\phi}$ as seen in the Klein-Gordon equation. The field is 
then expected to decay into matter and radiation, although the exact physical process behind that is unknown. The 
inflaton field is then treated as a perturbation around a mean field, that these perturbation give rise to perturbation amplitudes that we can identify in the observations we make today. 
Some useful parameters to define for a slow-roll inlfation are :
\begin{equation}
    \epsilon = - \frac{\dot{H}}{H^2} = \frac{3}{2} \frac{\dot{\phi}^2}{V}
\end{equation}
\begin{equation}
    \eta = \epsilon - \frac{\ddot{\phi}}{H \dot{\phi}}
\end{equation}
where the slow rolling conditions are $\epsilon \ll 1$ and $\eta \ll 1$. \\
We define $q$ as:
\begin{equation}
    q(x) = \delta \phi + \frac{\dot{\bar{\phi}}}{H} \Phi
\end{equation}
This quantity $q$ is known as the Mukhanov-Sasaki variable
(\cite{Mukhanov_1988}, \cite{Sasaki_1986}) and is the canonical variable for quantizing inflaton perturbations.\\

In order to fully treat the inflaton field we have to apply the canonical quantization of the field (quantum field theory) by elevating $q$ to an operator $\hat{q}$ as well as introducing 
its canonical conjugate momentum $\hat{\pi}$. In the Fourrier space, they are expressed as function of the quantum harmonic operator $a_{\mathbf{k}}$ and $a_{\mathbf{k}}^{\dag}$. 
With these operators we can calculate the expected value of the field as well as power spectrum as a function of $k$. 
Once the power spectrum during inflation is determined it is possible to extend this to the primordial power spectrum of matter during the matter dominated era. 
Outside the Hubble horizon, each mode is frozen with a nearly scale-invariant amplitude. 
Before writing the primordial power spectrum, we introduce the comoving curvature perturbation $\mathcal{R}$, a gauge-invariant 
quantity that measures the spatial curvature of constant-density hypersurfaces. In the context of inflation it is related to the inflaton perturbation $\delta\phi$ by:
\begin{equation}
    \mathcal{R} = -\frac{H}{\dot{\bar{\phi}}}\delta\phi
\end{equation}
The key property of $\mathcal{R}$ is that it is conserved on super-Hubble scales for adiabatic perturbations (\cite{Wands_2000}). 
This means its amplitude is frozen once a mode exits the Hubble radius during inflation and remains unchanged until it re-enters the horizon 
during the hot Big Bang, directly connecting inflationary predictions to late-time observables.
The dimensionless power spectrum of the comoving curvature perturbation 
$\mathcal{R}$ evaluated at Hubble crossing $k = aH$ is:
\begin{equation}
    \mathcal{P}_{\mathcal{R}}(k) = \frac{H^4}{4\pi^2 \dot{\bar{\phi}}^2}\bigg|_{k=aH} \approx \frac{H^2}{8\pi^2 M_{\mathrm{pl}}^2 \epsilon}\bigg|_{k=aH}
\end{equation}
which is approximately scale-invariant for slow-roll inflation where $H$ and $\epsilon$ vary slowly. The spectral index for a slow-roll inflation is given by:
\begin{equation}
    n_s = 1 - 6\epsilon + 2\eta
\end{equation}
Once modes re-enter the Hubble horizon during the matter and radiation dominated eras, they begin to evolve dynamically. Their evolution is captured by 
the transfer function $T(k)$ and the linear growth factor $D_+$. The linear matter power spectrum at a time $\tau$ is therefore:
\begin{equation}
    \mathcal{P}_{\mathrm{lin}}(\tau, k) = A_s \left(\frac{k}{k_*}\right)^{n_s} T^2(k)\, D_+^2(\tau)
\end{equation}
where $A_s$ is the primordial amplitude defined at the pivot scale $k_* = 0.05\,{\text{Mpc}}^{-1}$ (\cite{Planck_2018_inflation}), $T(k)$ encodes 
the suppression of small-scale modes that entered the horizon during radiation domination (\cite{Eisenstein_Hu_1998}), and $D_+(\tau)$ describes the  
growth of perturbations during matter domination (\cite{Carroll_1992}). The $k^{n_s}$ factor reflects the nearly scale-invariant primordial spectrum, with 
$n_s \approx 0.965$ (\cite{Planck_2018_inflation}) consistent with slow-roll inflation.

